% Part: normal-modal-logic
% Chapter: axioms-systems
% Section: introduction

\documentclass[../../../include/open-logic-section]{subfiles}

\begin{document}

\olfileid{nml}{axs}{int}

\olsection{प्रस्तावना}

मूलभूत मोडल भाषेसाठी मोडल प्रतिमानांद्वारे चिन्हार्थमीमांसा
आपल्याकडे आहे. !!a{formula} वैध असणे, म्हणजे सर्व
प्रतिमानांतील सर्व जगांत सत्य असणे, किंवा प्रतिमानांच्या
किंवा चौकटींच्या एखाद्या वर्गाच्या संदर्भात वैध असणे,
म्हणजे त्या वर्गातील अथवा त्या चौकटीवर आधारित सर्व
प्रतिमानांच्या सर्व जगांत सत्य असणे, या संकल्पनाही
आपल्याकडे आहेत. तर्कशास्त्रात अशा चिन्हार्थविषयक
वैधतेला सहसा !!{derivability} या सिद्धताविषयक
संकल्पनेशी जोडले जाते. एखाद्या पद्धतीत
!!{derivability} अशी परिभाषित करायची की
!!a{formula} वैध असेल तेव्हा आणि केवळ तेव्हाच
!!{derivable} असेल, हे आपले उद्दिष्ट आहे.

सर्वात सोप्या आणि ऐतिहासिकदृष्ट्या सर्वात जुन्या
!!{derivation} पद्धती म्हणजे हिल्बर्ट-प्रकारच्या किंवा
स्वयंसिद्धकीय !!{derivation} पद्धती. अनेक मोडल
तर्कशास्त्रांसाठी हिल्बर्ट-प्रकारच्या !!{derivation}
पद्धती तुलनेने सहज रचता येतात: अधितात्त्विक अभ्यासाच्या
दृष्टीने (उदा., निर्दोषता व संपूर्णता सिद्ध करताना)
त्या सोप्या असतात. पण त्यांमध्ये !!{formula}s सिद्ध
करणे, उदाहरणार्थ नैसर्गिक निगमन पद्धतींपेक्षा, बरेच
कठीण असते.

हिल्बर्ट-प्रकारच्या !!{derivation} पद्धतीत
!!a{formula} ची !!a{derivation} म्हणजे काही स्वयंसिद्धकांपासून
मोजक्या अनुमान नियमांनी मिळणारी, संबंधित
!!{formula} पर्यंत पोचणारी !!{formula}s ची मालिका.
!!{derivation} पद्धत चिन्हार्थमीमांसेशी जुळावी अशी
आपली अपेक्षा असल्याने, !!{derivable} सूत्रे
सर्व प्रतिमानांत (किंवा सर्व स्वयंसिद्धके सत्य असलेल्या
सर्व प्रतिमानांत) सत्य आहेत याची खात्री करावी लागेल.
यासाठी आवश्यक असलेले मोडल तर्कशास्त्रांचे काही
गुणधर्म आधी वेगळे पाहू: ``सामान्य'' मोडल तर्कशास्त्रे.
सामान्य मोडल तर्कशास्त्रांसाठी फक्त दोन अनुमान नियम
गृहीत धरावे लागतात: विध्यात्मक अनुमान आणि
अनिवार्यीकरण. स्वयंसिद्धके म्हणून सर्वतः-सत्य
सूत्रांचे सर्व आदेशन-दृष्टांत आणि संबंधित मोडल
तर्कशास्त्रानुसार काही मोडल स्वयंसिद्धके घेतो.
सर्व प्रतिमानांच्या वर्गातच आपल्याला रस असला,
तरी~$\Ax{K}$\iftag{notprvDiamond}{}{ आणि~$\Ax{Dual}$}
यांचे सर्व आदेशन-दृष्टांत स्वयंसिद्धके म्हणून घ्यावे
लागतात. एवढ्यानेच न्यूनतम सामान्य मोडल
तर्कशास्त्र~$\Log K$ मिळते.

\begin{defn}
\emph{विध्यात्मक अनुमान} हा नियम पुढील अनुमान-रूपबंध आहे:
\begin{prooftree}
\AxiomC{$!A$}
\AxiomC{$!A \lif !B$}
\RightLabel{\MP}
\BinaryInfC{$!B$}
\end{prooftree}
!!a{formula}~$!B$ हे !!{formula}s $!A$, $!C$
यांपासून विध्यात्मक अनुमानाने \emph{निष्पन्न होते}
तेव्हा आणि केवळ तेव्हाच, जेव्हा
$!C \ident !A \lif !B$.
\end{defn}

\begin{defn}
\emph{अनिवार्यीकरण} हा नियम पुढील अनुमान-रूपबंध आहे:
\begin{prooftree}
\AxiomC{$!A$}
\RightLabel{\Nec}
\UnaryInfC{$\Box !A$}
\end{prooftree}
!!{formula}~$!B$ हे !!{formula}s $!A$
पासून अनिवार्यीकरणाने निष्पन्न होते तेव्हा आणि
केवळ तेव्हाच, जेव्हा $!B \ident \Box !A$.
\end{defn}

\begin{defn}
स्वयंसिद्धकांच्या संच~$\Sigma$ पासूनची
\emph{!!{derivation}} म्हणजे !!{formula}s ची
$!B_1$, $!B_2$, \dots, $!B_n$ अशी मालिका,
ज्यात प्रत्येक $!B_i$ पुढीलपैकी एक असते:
\begin{enumerate}
\item सर्वतः-सत्य सूत्राचा आदेशन-दृष्टांत; किंवा
\item $\Sigma$ मधील !!a{formula} चा आदेशन-दृष्टांत; किंवा
\item $j$, $k < i$ असलेल्या दोन !!{formula}s
  $!B_j$, $!B_k$ यांपासून विध्यात्मक अनुमानाने
  निष्पन्न होणारे; किंवा
\item $j < i$ असलेल्या !!a{formula}~$!B_j$
  पासून अनिवार्यीकरणाने निष्पन्न होणारे.
\end{enumerate}
$!B_n \ident !A$ असलेली अशी !!a{derivation}
असल्यास $!A$ हे $\Sigma$ पासून
\emph{!!{derivable} आहे}, असे म्हणतो;
संकेतात $\Sigma \Proves !A$.
\end{defn}

या व्याख्येनुसार, !!{derivable} सूत्रांचा संच
सामान्य मोडल तर्कशास्त्र बनतो, आणि प्रत्येक
!!{derivable} !!{formula} हे सर्व स्वयंसिद्धके सत्य
असलेल्या प्रत्येक प्रतिमानात सत्य असते, हे पुढे
दिसेल. !!{derivation}s च्या या गुणधर्माला
\emph{निर्दोषता} म्हणतात. त्याचा व्यत्यास,
\emph{संपूर्णता}, सिद्ध करणे अधिक कठीण आहे.

\end{document}
